% Imporditud: tools/import_eksperimendid.py
% Allikas: Eesti füüsikaolümpiaadi eksperimendi ülesannete
%   kogu (Rael Kalda, 2023), ülesanne Ü114, lahendus L114.
\setAuthor{Oleg Košik}
\setRound{lõppvoor}
\setYear{2016}
\setNumber{P E1}
\setDifficulty{3}
\setTopic{elektriahelad}
\setVahendid{Patarei, tuntud takistusega takisti (R = 5,6 k$\Omega$), tundmatu takistusega takisti, voltmeeter, ühendusjuhtmed. Arvestage, et voltmeetril on olemas takistus}
\setPunktid{12}
\setAllikas{Eksperimendi kogu Ü114, lahendus lk 89}
\setLahendusseis{toores}

\prob{Takistid}
Määrake tundmatu takisti takistus.

\hint

\solu
Ühendame kõigepealt patarei voltmeetriga ja mõõdame

patarei pinge $U_0$ (näidiskatses saime $U_0$ = 4,9 V).

Nüüd leiame voltmeetri takistuse. Selleks ühendame

voltmeetri jadamisi patarei ja tuntud takistiga R =

5,6 k$\Omega$. Mõõdame voltmeetri pingeks $U_1$ (näidiskatses

$U_1$ = 2,2 V). Pinge takistil on seega UR = $U_0$ $-$ $U_1$, vool

$I_1$ = UR ning voltmeetri takistuseks saame R

RV = $U_1$ = U1R = U1R = 4,5 k$\Omega$. $I_1$ UR $U_0$ $-$ $U_1$

Lõpuks leiame tundmatu takisti takistuse r. Nüüd ühendame voltmeetri jadamisi patarei ja tundmatu takistiga. Mõõdame voltmeetri pingeks $U_2$ (näidiskatses $U_2$ = 3,1 V). Sarnaselt eelmise skeemiga leiame takisti takistuseks

r = ($U_0$ $-$ $U_2$)RV = 2,6 k$\Omega$. $U_2$
\probend
